%------------------------------------------------------------------------------------
%  SECTION 1 : PRÉSENTATION ET TYPOGRAPHIE
%------------------------------------------------------------------------------------
\section{Présentation et Typographie}
\label{sec:presentation}

Ce premier chapitre a pour objectif de démontrer les capacités graphiques et typographiques de ce modèle.
Chaque phrase est rédigée sur une ligne distincte afin de faciliter les revues de code et le versionnement sous Git.

\subsection{Familles de polices disponibles}
\label{subsec:polices}

Le modèle intègre plusieurs polices configurables :
\begin{itemize}
  \item \textbf{Police par défaut :} Calibri (sobre, fluide et parfaitement adaptée aux rapports académiques).
  \item \textbf{Police de titre / accentuation :} {\AsapFont Asap Condensed pour les titres principaux}.
  \item \textbf{Police spécifique :} {\ComicSansMSFont Comic Sans MS} (via \texttt{\textbackslash ComicSansMSFont}).
\end{itemize}

\subsection{Unités et formules scientifiques}
\label{subsec:sciences}

Le package \texttt{siunitx} permet d'écrire des grandeurs physiques sans ambiguïté :
Une vitesse de \qty{25}{\meter\per\second}, une température de \qty{37.2}{\degreeCelsius}, ou encore une masse de \qtyrange{50}{80}{\kilo\gram}.

Les équations mathématiques sont formatées selon les recommandations modernes :
\begin{equation}
  \label{eq:frequence_cardiaque}
  f_{\text{max}} = 220 - \text{âge} \pm \epsilon,
\end{equation}
où $\epsilon$ représente la marge d'erreur individuelle.
Une référence intelligente pointe vers l'\cref{eq:frequence_cardiaque} située à la \cpageref{eq:frequence_cardiaque}.

\subsection{Structuration par paragraphes avec puces en ligne}
\label{subsec:paragraphes}

La classe personnalise la commande \texttt{\textbackslash paragraph} pour insérer une puce discrète et enchaîner le texte directement sur la même ligne :

\paragraph{Critère de complétude}
Le document doit couvrir l'ensemble des axes d'analyse de manière exhaustive et rigoureuse.

\paragraph{Critère d'accessibilité visuelle}
Les contrastes de couleurs, la taille des polices et la structure hiérarchique doivent garantir un confort de lecture optimal.

\subsection{Encadrés et boîtes visuelles (tcolorbox)}
\label{subsec:boites}

Le package \texttt{tcolorbox} est configuré pour enrichir la mise en page d'éléments d'accroche :

\begin{boiteSynthese}{Synthèse méthodologique}
  Ce modèle applique les principes de modularité, de compilation reproductible et de séparation stricte entre le contenu et la mise en forme.
  Toutes les fonctionnalités additionnelles (annotations rétractables, encadrés, géométries adaptatives) sont prêtes à l'emploi.
\end{boiteSynthese}
